\documentclass{standalone}
\usepackage{circuitikz}
\usetikzlibrary{calc}
\definecolor{circuitnotemark}{HTML}{FE00FE}
\begin{document}
\begin{circuitikz}[american, line width=0.8pt]
\ctikzset{bipoles/length=1.2cm}
\foreach \x in {0,2,4,6,8,10,12,14,16} {\foreach \y in {0,-2,-4,-6,-8,-10,-12} {\fill[gray, opacity=0.35] (\x,\y) circle (1.1pt);}}
\node[gray, font=\scriptsize] at (0,0.84) {1};
\node[gray, font=\scriptsize] at (2,0.84) {2};
\node[gray, font=\scriptsize] at (4,0.84) {3};
\node[gray, font=\scriptsize] at (6,0.84) {4};
\node[gray, font=\scriptsize] at (8,0.84) {5};
\node[gray, font=\scriptsize] at (10,0.84) {6};
\node[gray, font=\scriptsize] at (12,0.84) {7};
\node[gray, font=\scriptsize] at (14,0.84) {8};
\node[gray, font=\scriptsize] at (16,0.84) {9};
\node[gray, font=\scriptsize, anchor=east] at (-0.84,0) {a};
\node[gray, font=\scriptsize, anchor=east] at (-0.84,-2) {b};
\node[gray, font=\scriptsize, anchor=east] at (-0.84,-4) {c};
\node[gray, font=\scriptsize, anchor=east] at (-0.84,-6) {d};
\node[gray, font=\scriptsize, anchor=east] at (-0.84,-8) {e};
\node[gray, font=\scriptsize, anchor=east] at (-0.84,-10) {f};
\node[gray, font=\scriptsize, anchor=east] at (-0.84,-12) {g};
\coordinate (a1) at (0,0);
\coordinate (a3) at (4,0);
\coordinate (e7) at (12,-8);
\coordinate (e9) at (16,-8);
\coordinate (g1) at (0,-12);
\coordinate (g3) at (4,-12);
\coordinate (c7) at (12,-4);
\coordinate (c9) at (16,-4);
\coordinate (a7) at (12,0);
\coordinate (c3) at (4,-4);
\draw (a1) to[R, l_=$R_{1}$] (a3); % line 2
\draw (e7) to[R, l_=$R_{2}$] (e9); % line 3
\draw (g1) to[R, l_=$R_{3}$] (g3); % line 4
\draw (c7) to[R, l_=$R_{4}$] (c9); % line 5
\draw (a3) -| (e7); % line 7
\draw (g3) |- (c7); % line 8
\node[circ] at (c7) {};
\path (20,-5.532) rectangle (24.024,0.395); % line 10
\node[anchor=west, inner sep=0, circuitnotemark, font=\footnotesize] at (20,0) {X}; % line 10
\node[anchor=west, inner sep=0, circuitnotemark, font=\footnotesize] at (20,-0.395) {X}; % line 10
\node[anchor=west, inner sep=0, circuitnotemark, font=\footnotesize] at (20,-0.79) {X}; % line 10
\node[anchor=west, inner sep=0, circuitnotemark, font=\footnotesize] at (20,-1.185) {X}; % line 10
\node[anchor=west, inner sep=0, circuitnotemark, font=\footnotesize] at (20,-1.58) {X}; % line 10
\node[anchor=west, inner sep=0, circuitnotemark, font=\footnotesize] at (20,-1.976) {X}; % line 10
\node[anchor=west, inner sep=0, circuitnotemark, font=\footnotesize] at (20,-2.371) {X}; % line 10
\node[anchor=west, inner sep=0, circuitnotemark, font=\footnotesize] at (20,-2.766) {X}; % line 10
\node[anchor=west, inner sep=0, circuitnotemark, font=\footnotesize] at (20,-3.161) {X}; % line 10
\node[anchor=west, inner sep=0, circuitnotemark, font=\footnotesize] at (20,-3.556) {X}; % line 10
\node[anchor=west, inner sep=0, circuitnotemark, font=\footnotesize] at (20,-3.951) {X}; % line 10
\node[anchor=west, inner sep=0, circuitnotemark, font=\footnotesize] at (20,-4.346) {X}; % line 10
\node[anchor=west, inner sep=0, circuitnotemark, font=\footnotesize] at (20,-4.741) {X}; % line 10
\node[anchor=west, inner sep=0, circuitnotemark, font=\footnotesize] at (20,-5.136) {X}; % line 10
\end{circuitikz}
\end{document}
